\documentclass[twocolumn]{iopjournal}

\usepackage{amsmath,amssymb,amsfonts}
\usepackage{amsthm}

\fancyhfoffset[L]{0mm}
\setlength{\marginparwidth}{0pt}
\setlength{\marginparsep}{0pt}

% ==========================================
% ZENODO PREPRINT HEADERS (ACTIVE)
% ==========================================
\renewcommand{\articletype}[1]{%
  {\vspace*{-8mm}\noindent \Large \sf Preprint}\par
  \vspace*{6mm}\noindent{\scriptsize \sf{\bfseries \MakeUppercase{#1}}}}

\fancyhead[L]{}
\fancyhead[R]{}

\hypersetup{pdftitle={The Constitutive Law as a Tug of War}, pdfauthor={Stephen Nelson}, pdfsubject={Gravity Field Dynamics (GFD)}, pdfkeywords={constitutive law, interpolating function, MOND, exchange symmetry, screening, decoding mass, vortex, proof}}

\newtheorem{lemma}{Lemma}
\newtheorem{theorem}{Theorem}
\newtheorem{proposition}{Proposition}
\newtheorem{corollary}{Corollary}
\theoremstyle{definition}
\newtheorem{definition}{Definition}
\theoremstyle{remark}
\newtheorem*{remark}{Remark}

\begin{document}

\articletype{Proof}

\title{The Constitutive Law as a Tug of War}

\author{Stephen Nelson}

\affil{Field Dynamics Research}

\ead{stephen@fielddynamics.org}

\begin{abstract}
In Gravity Field Dynamics the constitutive law $\mu(x) = x/(x+1)$, with
$x = g/a_0$, is the share of the field carried by TetA. We show that it is
the outcome of a tug of war in which TetA pulls with the local field $g$
and TetB pulls with $a_0$, so that $\mu = g/(g+a_0)$, and we identify
exactly which premises make it so. On the cube, every state with the
tetrahedral symmetry holds one value on each chamber, and the vortex modes
cannot move the share. When a law obeys a screening relation of
mass-action form, $x^n(1-\mu)^m = \mu^p$, the exchange of the two chambers
holds exactly when $m = p$, and its odds $\mu/(1-\mu)$ are then a power of
$x$. Among those powers, the first is singled out by each of three
equivalent conditions: the two strengths are accelerations, rotation
curves are flat, and the relation has degree one. Every other premise is
shown necessary by an explicit alternative. The same algebra reads the
decoding-mass letter's drive and brake as TetB pulling harder or softer
than $a_0$: the decoded mass exceeds the photometric mass by
$(w - a_0)/(g + a_0)$, where $w$ is TetB's measured strength. Exact
Navier--Stokes solutions supply the vortex picture behind that reading.
\end{abstract}

\keywords{constitutive law, interpolating function, MOND, exchange
symmetry, screening, decoding mass, Burgers vortex, proof}

% ======================================================================

\section*{Setting}

At a point of a galaxy's net field the total pull is $g$, the pull of the
visible matter is $g_{\rm bar}$, and $x = g/a_0$, with $a_0$ Milgrom's
acceleration scale~\cite{Milgrom1983}. AQUAL~\cite{BM1984}
relates them through $\mu(x)\,g = g_{\rm bar}$, and in GFD $\mu$ is the
share of the field carried by TetA, the chamber we observe, while $1 - \mu$
is carried by TetB, which lies beyond the Field Origin and is seen only
through its share~\cite{NelsonL2,NelsonL3}. A share is most naturally
written as odds,
\begin{equation}
  \Omega = \frac{\mu}{1-\mu},
\end{equation}
which run from $0$ (all TetB) to $\infty$ (all TetA) and equal $1$ at an
even split. The exchange of the two chambers sends $\Omega$ to $1/\Omega$,
and it is written
\begin{equation}
  \mu(1/x) = 1 - \mu(x). \tag{X}
\end{equation}
Taken at $x = 1$, (X) gives $\mu(1) = 1/2$, so the scale $a_0$ marks the
pull at which the two chambers carry equal shares.

The question this proof answers is the one a referee asks first: AQUAL
admits infinitely many interpolating functions with the right limits, so
what, exactly, selects $x/(x+1)$? Earlier routes~\cite{NelsonCL,NelsonCLL}
each identify the odds with the acceleration ratio; here that
identification is decomposed into premises that can be tested one at a
time.

\section*{Two teams on the cube}

\begin{lemma}[Two teams]\label{lem:teams}
Let a state on the cube's eight corners have the tetrahedral symmetry
$T_d$. Then it holds one value $w_A$ on TetA's four corners and one value
$w_B$ on TetB's, and
\begin{enumerate}
\item[(i)] its TetA share is $\mu = w_A/(w_A + w_B)$;
\item[(ii)] the ratio of its $xyz$ and constant Walsh coefficients is
  $s = (w_A - w_B)/(w_A + w_B) = 2\mu - 1$;
\item[(iii)] the inversion $-I$ swaps $w_A$ and $w_B$, so $\mu \to 1-\mu$;
\item[(iv)] adding any combination of the six $T_2$ modes leaves $\mu$
  unchanged.
\end{enumerate}
\end{lemma}

\begin{proof}
$T_d$ acts transitively on each chamber, so a symmetric state is constant
on each, and (i)--(iii) are then direct sums over the corners. The $T_2$
modes sum to zero over each chamber, which is (iv). All four are proved in
the Walsh-eigenmodes formalisation~\cite{NelsonWalsh}.
\end{proof}

The lemma says that the geometry leaves exactly two numbers free, one
strength per team, and that their competition is all the share can
express. The quantity $s$ is the position of the rope: it is the balance
point of a corner $v$ weighted by $w_A$ and its antipode $-v$ weighted by
$w_B$ on the body diagonal through the Field Origin, and with $w_A = x$,
$w_B = 1$ it is the Möbius map $s = (x-1)/(x+1)$ of the archive's coupling
polynomial, whose capacity $f = 17 + 4s = (21x+13)/(x+1)$ follows
at once~\cite{NelsonDocs}. Item (iv) matters for what follows, because the
cube's vortices are built from the $T_2$ modes~\cite{NelsonWalsh}: whatever
rotation does, it cannot move the rope.

\section*{Exchange turns screening into a power law}

The constitutive-law synthesis~\cite{NelsonSynth} writes the law as a
screening balance, signal times remaining capacity equals used capacity,
$x(1-\mu) = \mu$. Its natural generalisation is the mass-action form
\begin{equation}
  x^n (1-\mu)^m = \mu^p, \qquad n, p > 0, \tag{S}
\end{equation}
which contains the synthesis law $(1,1,1)$, its bimetric law
$x(1-\mu)^2 = \mu$, the bound-state law of Paper~XIII, and the Hill forms of
cooperative binding~\cite{Hill1910}.

\begin{theorem}\label{thm:power}
Let $0 < \mu(x) < 1$ on $x > 0$ obey (S). Then $\mu$ satisfies the
exchange (X) if and only if $m = p$. When $m = p$ the odds are a pure
power,
\begin{equation}
  \Omega(x) = x^k, \qquad \mu = \mu_k(x) = \frac{x^k}{1 + x^k}, \qquad
  k = \frac{n}{p},
\end{equation}
and every $\mu_k$ satisfies (X).
\end{theorem}

\begin{proof}
If (X) holds, then $\mu(1) = 1/2$, and (S) at $x = 1$ reads
$(1/2)^m = (1/2)^p$, so $m = p$. Conversely, if $m = p$, (S) says
$\Omega^p = x^n$, so $\Omega = x^{n/p}$; and
$\mu_k(1/x) = 1/(1+x^k) = 1 - \mu_k(x)$.
\end{proof}

The theorem converts a symmetry into a shape. The exchange, a statement
about one reflection, removes every law in which the screening powers
differ, and what survives is a family with a single exponent. Scale
covariance, $\Omega(\lambda x) = \Omega(\lambda)\,\Omega(x)$, which an
earlier draft of this argument had to assume, now follows.

\section*{What fixes the power}

\begin{theorem}\label{thm:k}
For the family $\mu_k$, the following are equivalent:
\begin{enumerate}
\item[(i)] $k = 1$;
\item[(ii)] the two strengths are accelerations that add up to the field:
  $\mu_k(g/a_0)\,(g + a_0) = g$ for every $g > 0$;
\item[(iii)] rotation curves are flat in the deep field;
\item[(iv)] the odds are a degree-one (Möbius) function of $x$.
\end{enumerate}
When they hold, $\mu = g/(g+a_0) = x/(x+1)$, TetB's strength is $a_0$, and
$v^4 = GMa_0$.
\end{theorem}

\begin{proof}
(i)$\Leftrightarrow$(ii): $\mu_k(x)(x+1) = x$ is equivalent to $x^k = x$,
and at $x = 2$ that forces $k = 1$. (i)$\Leftrightarrow$(iii): in the deep
field $\mu_k \simeq x^k$, so $g^{k+1} = GMa_0^k/r^2$ and the circular speed
obeys
\begin{equation}
  v^2 = g\,r = (GMa_0^k)^{1/(k+1)}\, r^{(k-1)/(k+1)},
\end{equation}
which is independent of $r$ exactly when $k = 1$, and then $v^4 = GMa_0$.
(i)$\Leftrightarrow$(iv): $x^k$ is Möbius only for $k = \pm 1$, and $k > 0$.
\end{proof}

Condition (ii) is the tug of war itself. With $\Omega = x^k$ the strengths
are $g^k$ and $a_0^k$, and only for $k = 1$ are they accelerations whose
sum, weighted by the share, returns the field. So the observation that
$a_0$ is an acceleration, on the same footing as $g$, is what fixes the
exponent, and flat rotation curves fix it independently. Condition (iv) is
the premise (A) of the exchange-symmetry proof~\cite{NelsonCL}, which
reaches the same law from degree one, the endpoint $\mu(0) = 0$ and
closure under exchange; the two routes agree because they meet only at
$k = 1$.

\begin{proposition}\label{prop:need}
Each premise is necessary.
\begin{enumerate}
\item[(a)] Without (S), a law can keep the exchange, both MOND limits
  and monotonicity, and still differ from $x/(x+1)$: for $c \ne 0$ the
  odds
  \[\Omega = x\exp\bigl(c\,t\,e^{-t^2}\bigr),\qquad t = \ln x,\]
  do all of this.
\item[(b)] Without $m = p$, the exchange fails: $x(1-\mu)^2 = \mu$ has
  $\mu(1) = (3-\sqrt5)/2 = 1/\varphi^2$, not $1/2$.
\item[(c)] Without $k = 1$, rotation curves are not flat: $k = 2$ gives
  $v^2 \propto r^{1/3}$.
\item[(d)] The standard function $x/\sqrt{1+x^2}$ and the exponential
  function $1 - e^{-x}$ break the exchange.
\item[(e)] Degree one, $\mu(0) = 0$ and the exchange leave two branches,
  $x/(x+1)$ and $x/(x-1)$, and only the first stays in $[0,1]$.
\end{enumerate}
\end{proposition}

\begin{proof}
(a) $\ln\Omega - t$ is odd in $t$, which is (X), and vanishes as
$t \to \pm\infty$, which gives both limits. (b) and (d) follow because (X)
forces $\mu(1) = 1/2$, while these laws give $1/\varphi^2$, $1/\sqrt2$ and
$1 - e^{-1}$. (c) is Theorem~\ref{thm:k}. (e) is the coefficient
comparison of the degree-one relation.
\end{proof}

Item (b) is worth stating plainly. The bound-state law is not wrong; it
describes a self-sourcing field, which couples in both directions, and the
proposition shows that such a field cannot be exchange-symmetric. The
11.8\% gap between $1/2$ and $1/\varphi^2$ at $x = 1$, which the
synthesis~\cite{NelsonSynth} proposes as a test with SPARC~\cite{SPARC}, is
therefore a test of the exchange itself.

\section*{The Field Origin and the cap}

\begin{proposition}\label{prop:node}
Between the two teams the following hold.
\begin{enumerate}
\item[(i)] A single node between two teams, in balance
  $c_A(V_A - V_O) + c_B(V_B - V_O) = 0$, sits at their weighted average
  $V_O = (c_AV_A + c_BV_B)/(c_A + c_B)$, and eliminating it leaves the
  series link $c_Ac_B/(c_A+c_B)$.
\item[(ii)] Each team carries $\mu$ times its own strength:
  $g = \mu g + \mu a_0$. TetB therefore carries $(1-\mu)g = \mu a_0$, which
  rises with $g$ towards $a_0$ and never reaches it, and the missing pull
  over the visible pull is $a_0/g$.
\item[(iii)] A swap of strong and weak pulls that commutes with rescaling
  is $g \mapsto c/g$, and its only fixed pull is $\sqrt c$, which is $a_0$.
\end{enumerate}
\end{proposition}

\begin{proof}
Each item is a line of algebra; for (iii), set the scale equal to $g$ in
$E(\lambda g) = E(g)/\lambda$.
\end{proof}

Item (i) is why the Field Origin matters. If the two chambers meet only
there, the centre is a single node between the teams; eliminating it from
the cube, the Schur complement proved in~\cite{NelsonWalsh}, leaves a
single rank-one rope whose TetA share is $t_A/(t_A + t_B)$, and direct
links along the cube's edges would break that rule. Item (ii) explains the
statement of Letters~2 and~3 that the TetB field never exceeds
$a_0$~\cite{NelsonL2,NelsonL3}: TetB's pull is its share of its own
strength.

\section*{Decoding mass: quiet, drive and brake}

The decoding-mass letter~\cite{NelsonDecode} turns an observed rotation
curve into $M_{\rm dec} = rv^4/[G(a_0r + v^2)]$, which is
$\mu(g/a_0)\,g\,r^2/G$ with $g = v^2/r$, and compares it with the
photometric mass $M_{\rm phot} = g_{\rm bar} r^2/G$. The measured share is
$g_{\rm bar}/g$, and the strength that TetB would need for that share to be
the quiet one is
\begin{equation}
  w = \frac{g\,(g - g_{\rm bar})}{g_{\rm bar}}.
\end{equation}

\begin{theorem}\label{thm:decode}
At every radius,
\begin{equation}
  \frac{M_{\rm dec}}{M_{\rm phot}} = \frac{g + w}{g + a_0},\qquad
  \frac{\Delta M}{M_{\rm phot}} = \frac{w - a_0}{g + a_0}.
\end{equation}
The decode is quiet, $M_{\rm dec} = M_{\rm phot}$, exactly when $w = a_0$,
that is, when $g_{\rm bar} = \mu g$; it shows drive,
$M_{\rm dec} > M_{\rm phot}$, exactly when $w > a_0$, that is, when TetB
pulls more than $\mu a_0$; and it shows brake exactly when $w < a_0$.
\end{theorem}

\begin{proof}
Substitute $\mu = g/(g+a_0)$; then
$g - g_{\rm bar} - \mu a_0 = (g_{\rm bar}/g)\,\mu\,(w - a_0)$, whose prefactor
is positive.
\end{proof}

The theorem gives $a_0$ a second meaning: it is TetB's quiet strength.
Drive and brake are TetB, in the interior where the framework places
them~\cite{NelsonHourglass}, pulling harder or softer than that, and the
decoded mass is how the difference is read. The denominator $g + a_0$
explains where the reading works. The same change in TetB's strength moves
the decode by tens of per cent in a galaxy's outskirts and by parts in
$10^7$ where $g \gg a_0$, which is why drive and brake never appear in the
solar system.

\section*{The vortex}

\begin{proposition}\label{prop:vortex}
Three exact solutions of the Navier--Stokes equations describe a drained
vortex.
\begin{enumerate}
\item[(i)] (Burgers~\cite{Burgers1948}.) With inflow $u_r = -\alpha r/2$
  and axial outflow $u_z = \alpha z$, the vorticity
  $\omega = \omega_0\exp(-\alpha r^2/4\nu)$ is steady: inward advection,
  stretching and diffusion balance exactly,
  $u_r\,\omega' = \alpha\omega + \nu(\omega'' + \omega'/r)$.
\item[(ii)] (Lamb--Oseen.) With no drain, the same circulation diffuses:
  $\omega = (\Gamma/4\pi\nu t)\exp(-r^2/4\nu t)$ solves
  $\partial_t\omega = \nu(\omega'' + \omega'/r)$.
\item[(iii)] At fixed circulation $\Gamma$, the Burgers core has radius
  proportional to $\sqrt{\nu/\alpha}$ and central vorticity
  $\Gamma\alpha/4\pi\nu$; far from the core the swirl tends to
  $\Gamma/2\pi r$, whatever $\alpha$ is.
\end{enumerate}
\end{proposition}

\begin{proof}
Direct differentiation, and $\int_0^\infty e^{-\beta r^2}2\pi r\,dr =
\pi/\beta$.
\end{proof}

A drained vortex is sustained by its drain: in a bathtub the swirl is
maintained by gravity's work on the water falling through it, and its
spin-up comes from stretching the rotating column~\cite{Lautrup}. The
proposition makes the three regimes exact. When the drain balances
diffusion the vortex is steady, which is the quiet state; a stronger drain
tightens the core and spins it up, which is drive; a weaker drain lets it
spread, and with no drain at all it decays, which is brake carried to its
end. Item (iii) adds the observation that matches Theorem~\ref{thm:decode}:
the drain's strength is written in the core and in the inflow, not in the
far swirl, just as drive and brake show where TetB's share is large and
hide where $g \gg a_0$. The identification of TetB with the drain is the
framework's reading~\cite{NelsonHourglass}; the vortex equations and the
decode identity are exact.

\section*{Scope}

Lemma~\ref{lem:teams} and the exchange as an operation are geometry. Three
premises are physical: the screening form (S), the reading of the chamber
swap as $x \to 1/x$, and the statement that the strengths are
accelerations, which flat rotation curves confirm independently. The value
of $a_0$ is measured. Two statements of the March 2026
synthesis~\cite{NelsonSynth} need rewording without changing the law: $N$
independent sites fill to $x/(1+x)$ for every $N$, so the Hill form
$x^N/(1+x^N)$ requires cooperative filling~\cite{Hill1910,Langmuir1918};
and the inverse of $x/\sqrt{1+x^2}$ is algebraic but not rational.

\section*{Verification}

The script \texttt{verify\_constitutive\_law\_tug\_of\_war.py} runs 42
checks in eight parts: exactly with sympy where a statement is algebraic,
exhaustively over the cube's corners and symmetries, and numerically for
the rotation curves. The test file \texttt{test\_tug\_of\_war.py} in the
research folder \texttt{constitutive-law-u(x)} holds 85 finer tests of the
same statements.

% Lean: new paragraph.
\textbf{Formal verification.} Lemma~\ref{lem:teams},
Theorem~\ref{thm:power}, Theorem~\ref{thm:k} (i)--(iii),
Proposition~\ref{prop:need} (b)--(d), Proposition~\ref{prop:node},
Theorem~\ref{thm:decode} and Proposition~\ref{prop:vortex} (i)--(ii) are
checked in Lean~4 with Mathlib (\texttt{lean/TugOfWar.lean}), importing the
Walsh-eigenmodes and constitutive-law formalisations. Item (iv) of
Theorem~\ref{thm:k} is the constitutive-law formalisation's
\texttt{unique\_law}; the counterexample (a), the branch count (e) and the
Burgers core scaling are checked by the script.

\begin{thebibliography}{99}

\bibitem{BM1984}
J.~Bekenstein and M.~Milgrom,
Does the missing mass problem signal the breakdown of Newtonian
gravity?, \textit{Astrophys. J.} \textbf{286}, 7 (1984).

\bibitem{Milgrom1983}
M.~Milgrom,
A modification of the Newtonian dynamics as a possible alternative to
the hidden mass hypothesis, \textit{Astrophys. J.} \textbf{270}, 365
(1983).

\bibitem{NelsonCL}
S.~Nelson,
\textit{The Constitutive Law from Exchange Symmetry},
Field Dynamics Research (2026),
\url{proofs/constitutive-law-exchange-symmetry}.

\bibitem{NelsonCLL}
S.~Nelson,
\textit{The Gravitational Constitutive Law from the Tetrad Frame
Geometry}, Field Dynamics Research (2026).

\bibitem{NelsonL2}
S.~Nelson,
\textit{The Other Half of Gravity: Torsion and the Two-Chamber Structure
of the Tetrad Frame} (Letter~2),
Field Dynamics Research (2026).

\bibitem{NelsonL3}
S.~Nelson,
\textit{Gravity Field Dynamics: Field Closure of the Tetrad Frame}
(Letter~3), Field Dynamics Research (2026).

\bibitem{NelsonWalsh}
S.~Nelson,
\textit{The Walsh Eigenmodes of the Cube, and What Field Closure Reads},
Field Dynamics Research (2026),
\url{proofs/walsh-eigenmodes-of-the-cube}.

\bibitem{NelsonDocs}
S.~Nelson,
\textit{Constitutive Law $\mu(x)$}, documentation, Field Dynamics
Research (2026), \url{docs/02-gravitation}.

\bibitem{NelsonSynth}
S.~Nelson,
\textit{The Constitutive Law of Gravity Field Dynamics}, synthesis and
validation suites, Field Dynamics Research (2026),
\url{research/constitutive-law-u(x)}.

\bibitem{NelsonDecode}
S.~Nelson,
\textit{Gravity Field Dynamics: Decoding Mass from Observed Rotation
Curves},
Field Dynamics Research (2026).

\bibitem{NelsonHourglass}
S.~Nelson,
\textit{TetA and TetB as One Vortex}, research note, Field Dynamics
Research (2026), \url{research/vortex-hourglass}.

\bibitem{SPARC}
F.~Lelli, S.~S.~McGaugh and J.~M.~Schombert,
SPARC: mass models for 175 disk galaxies with Spitzer photometry at
3.6\,$\mu$m and accurate rotation curves, \textit{Astron. J.}
\textbf{152}, 157 (2016).

\bibitem{Burgers1948}
J.~M.~Burgers,
A mathematical model illustrating the theory of turbulence,
\textit{Adv. Appl. Mech.} \textbf{1}, 171 (1948).

\bibitem{Lautrup}
B.~Lautrup,
\textit{Physics of Continuous Matter}, 2nd ed. (CRC Press, Boca Raton,
2011), ch.~19.

\bibitem{Hill1910}
A.~V.~Hill,
The possible effects of the aggregation of the molecules of
haemoglobin on its dissociation curves, \textit{J. Physiol.}
\textbf{40}, iv (1910).

\bibitem{Langmuir1918}
I.~Langmuir,
The adsorption of gases on plane surfaces of glass, mica and platinum,
\textit{J. Am. Chem. Soc.} \textbf{40}, 1361 (1918).

\end{thebibliography}

\end{document}
